\chapter{Core Concepts}
\label{ch:core-concepts}

\newthought{A handful of ideas} recur on every page of this manual: the task card, the
standalone project, the path markers and runtime tokens that make paths portable, and the
fixed contract for where output lands. Understand these five and the rest of the book is
detail. None of them require code---they are conventions \Jarvis\ enforces so that scans stay
self-contained and reproducible.

\section{The task card}
\label{sec:task-card}

A \emph{task card} is a single \textsc{yaml} file that fully describes one scan. It is the only
thing you normally edit. As Chapter~\ref{ch:introduction} put it, every card fills in the same
three steps for each parameter point: draw a point, run a workflow to get observables, and
score them with a \yamlkey{LogLikelihood}.

A card is organised into a small set of top-level blocks. You will meet each in detail in
Part~\ref{part:task-card}; the shape is:

\begin{yamlwin}
Scan:         # run name + where outputs go
Sampling:     # how points are drawn + the objective
Calculators:  # external-program workflow   (choose one backend)
Operas:       # in-process workflow         (choose one backend)
EnvReqs:      # platform + dependency contract
LibDeps:      # shared external libraries (optional)
Utils:        # helper functions (optional)
\end{yamlwin}

\noindent \yamlkey{Scan} and \yamlkey{Sampling} are required; \yamlkey{EnvReqs} is strongly
recommended; the rest are optional, with one important rule covered in
Section~\ref{sec:backends}.

\section{Standalone projects}
\label{sec:standalone-projects}

\Jarvis\ deliberately has \emph{no repo-root requirement}. You do not run it from a checkout of
the framework; you run it from \emph{your} project directory. A project is just a folder
containing two \emph{marker files}:

\begin{description}[style=nextline, leftmargin=1.2em]
  \item[\file{.jarvis-project.json}] machine-readable project metadata;
  \item[\file{jarvis.project.yaml}] human-editable project metadata.
\end{description}

\noindent When \Jarvis\ starts, it walks upward from your task card looking for either marker
and treats the directory that holds it as the \emph{task root}. Everything project-local is
resolved relative to that root. \cli{Jarvis project create} (Chapter~\ref{ch:quickstart})
writes both markers for you.

\begin{jnote}
You can override the task root explicitly by setting the environment variable
\code{JARVIS\_HEP\_TASK\_ROOT} (or \code{JHEP\_TASK\_ROOT}). If no marker is found and no
override is set, \Jarvis\ falls back to the current working directory. Setting the markers is
the robust, recommended way.
\end{jnote}

\section{Path markers}
\label{sec:path-markers}

Hard-coded absolute paths make a project impossible to move or share. \Jarvis\ solves this with
a \emph{path marker}: a short symbol that expands to a real directory at runtime.

The one public path marker is \marker{\&J}:

\begin{table}[ht]
  \footnotesize
  \begin{tabular}{@{}ll@{}}
    \toprule
    \textbf{Marker} & \textbf{Expands to} \\
    \midrule
    \marker{\&J/...} & the standalone project root (the task root of
                       Section~\ref{sec:standalone-projects}) \\
    \bottomrule
  \end{tabular}
  \caption[The public path marker.]{The public path marker. Always write project-local paths
           with \marker{\&J/}.}
  \label{tab:markers}
\end{table}

\noindent So \marker{\&J/outputs} means ``the \file{outputs/} directory inside this project,''
wherever the project happens to live. Because the marker is resolved fresh on every run, the
same card works on your laptop, a collaborator's machine, or an \textsc{hpc} node with no
edits.

\begin{jnote}
\marker{\&J/} is the only path marker you need. Keep everything a scan depends on inside the
project and reference it with \marker{\&J/...}; for a shared default file, keep a copy under
\marker{\&J/deps/...} and point to that---exactly what the built-in quick-starts do.
\end{jnote}

\section{Runtime tokens}
\label{sec:runtime-tokens}

Where a path marker names a \emph{directory}, a \emph{runtime token} names a value that changes
\emph{per sample}. Tokens let one calculator definition give every sample its own isolated
working directory and file names. There are three (Table~\ref{tab:tokens}).

\begin{table}[ht]
  \footnotesize
  \begin{tabular}{@{}lp{0.62\linewidth}@{}}
    \toprule
    \textbf{Token} & \textbf{Substituted with} \\
    \midrule
    \token{@SampleID} & the current sample's \textsc{uuid} \\
    \token{@Sdir}     & the current sample's save directory, under
                        \file{outputs/.../SAMPLE/<uuid>} (or \file{.../SAMPLE/tests/<uuid>}
                        under \cli{-{}-check-modules}) \\
    \token{@PackID}   & a calculator module \emph{instance} id, enabling per-instance working
                        directories so parallel samples never clobber each other \\
    \bottomrule
  \end{tabular}
  \caption[Runtime tokens.]{Runtime tokens, substituted per sample on calculator workflow
           paths.}
  \label{tab:tokens}
\end{table}

\noindent These tokens are available on calculator workflow paths---commands, working
directories, and sample-scoped input/output file paths. You will use \token{@PackID} heavily
when wrapping external programs (Chapter~\ref{ch:calculators}); \token{@Sdir} when you want a
calculator to write into the per-sample artifact directory.

\begin{jnote}
Referencing \token{@Sdir} (or running any calculator workflow) causes \Jarvis\ to
\emph{materialize} a per-sample directory on disk. Pure-\yamlkey{Operas} scans can skip that
filesystem work entirely---a performance lever discussed in
Chapter~\ref{ch:runtime-block}.
\end{jnote}

\section{The output contract}
\label{sec:output-contract}

Every run writes a predictable tree, derived from \yamlkey{Scan.name} (written \code{<scan>}
below). Knowing this layout means you always know where to look.

\begin{table}[ht]
  \footnotesize
  \begin{tabular}{@{}p{0.45\linewidth}p{0.47\linewidth}@{}}
    \toprule
    \textbf{Location} & \textbf{Contents} \\
    \midrule
    \file{outputs/<scan>/DATABASE/} & \textsc{hdf5} samples, schema, \textsc{csv} exports, run metadata \\
    \file{outputs/<scan>/SAMPLE/}   & per-sample working files and retained artifacts \\
    \file{logs/<scan>/}             & Jarvis, sampler, and runtime logs \\
    \file{images/<scan>/}           & semantic \file{flowchart.json}, plot configs, figures \\
    \file{outputs/<scan>/run\_summary.\{json,csv,txt\}} & end-of-run summaries \\
    \bottomrule
  \end{tabular}
  \caption[The output contract.]{The project-local output contract. Every run follows it.}
  \label{tab:output-contract}
\end{table}

\noindent Two details worth remembering now:

\begin{itemize}
  \item The \file{DATABASE/} path is shared between a normal run and a \cli{-{}-check-modules}
        dry run; only the \file{SAMPLE/} side differs (the check writes under
        \file{SAMPLE/tests/}).
  \item \Jarvis\ always writes \file{flowchart.json} (a semantic description of your workflow);
        if \JPlot\ is installed it additionally renders \file{flowchart.png}. The
        \cli{-{}-skip-draw-flowchart} flag suppresses only the \textsc{png} step.
\end{itemize}

Part~\ref{part:running} returns to each of these artifacts in depth.

\section{Choosing a workflow backend}
\label{sec:backends}

The workflow---step two of the loop---is provided by \emph{one} of two backends:

\begin{description}[style=nextline, leftmargin=1.2em]
  \item[\yamlkey{Calculators}] drives external programs through files and shell commands. Use
        it for compiled \textsc{hep} packages and file-based toolchains
        (Chapter~\ref{ch:calculators}).
  \item[\yamlkey{Operas}] evaluates in-process Python operators. Use it for algebraic or
        array-based work, and for fast prototyping (Chapter~\ref{ch:operas}).
\end{description}

\begin{jwarn}
A task card uses \emph{either} \yamlkey{Calculators} \emph{or} \yamlkey{Operas} as its workflow
backend, not both. Pick the one that matches your tools; the quick-starts in
Chapter~\ref{ch:quickstart} use \yamlkey{Operas} because the \code{EggBox} toy is an in-process
operator.
\end{jwarn}

\section{Putting it together}
\label{sec:concepts-together}

With the vocabulary in hand, the whole framework reads cleanly: you author a \emph{task card}
inside a \emph{standalone project}; \marker{\&J/} keeps its paths portable and runtime tokens
isolate each sample; a sampler draws points, your chosen \emph{backend} turns them into
observables, a \yamlkey{LogLikelihood} scores them, and the results land in the fixed
\emph{output contract}. The next part dissects the task card block by block.
